\documentclass[11pt]{article}
\usepackage[margin=1in]{geometry}
\usepackage[T1]{fontenc}
\usepackage{amsmath,amssymb,amsthm}
\usepackage[hidelinks]{hyperref}

\newtheorem*{theorem}{Literature-Implied Answer}

\title{No Computable Uniform Dimension Bound for Open Problem 5.3 of arXiv:1512.00419}
\author{FA/Banach run packet}
\date{2026-06-29}

\begin{document}
\maketitle

\section*{Status}

This is a \texttt{literature\_implied\_answers} packet.  It records that
Open Problem 5.3 of Palazuelos--Vidick \cite{PVSurvey}, interpreted as asking
for an explicit computable dimension bound, has a negative answer as a
consequence of the later theorem \( \mathrm{MIP}^{*}=\mathrm{RE} \) of
Ji--Natarajan--Vidick--Wright--Yuen \cite{MIPRE}.  This is a status
identification and reduction, not a new nonlocal-game construction.

\section*{Source Question}

In Section 5.4.1, after asking whether finite games may require
infinite-dimensional entanglement, the source paper asks:

\begin{quote}
Given \(\varepsilon>0\) and \(n\in\mathbb N\), provide an explicit bound
\(d=d(\varepsilon,n)\) such that for any game \(G\) with at most \(n\)
questions and answers per player there exists a \(d\)-dimensional strategy
with success probability at least \(\omega^*(G)-\varepsilon\).
\end{quote}

The source notes that compactness and discretization give existence of some
finite \(d(\varepsilon,n)\), but not an estimate.

\section*{Answer}

\begin{theorem}
There is no total computable function \(D(\varepsilon,n)\), for positive
rational \(\varepsilon\) and \(n\in\mathbb N\), such that every finite
two-player nonlocal game with at most \(n\) questions and answers per player
has a local-dimension-\(D(\varepsilon,n)\) strategy with value at least
\(\omega^*(G)-\varepsilon\).
\end{theorem}

\begin{proof}
Ji--Natarajan--Vidick--Wright--Yuen prove
\(\mathrm{MIP}^{*}=\mathrm{RE}\).  Their abstract states, in particular, that
there is an efficient reduction from the Halting Problem to deciding whether a
two-player nonlocal game has entangled value \(1\) or at most \(1/2\)
\cite{MIPRE}.

Suppose, toward a contradiction, that such a computable bound \(D\) exists.
Given a game \(G\) produced by the reduction, compute a natural number \(n\)
bounding its question and answer alphabets and set
\[
  d = D(1/6,n).
\]
If \(\omega^*(G)=1\), then by the assumed dimension bound there is a
\(d\)-dimensional strategy with success probability at least \(5/6\).  If
\(\omega^*(G)\le 1/2\), then every strategy, in particular every
\(d\)-dimensional strategy, has success probability at most \(1/2\).

For fixed \(G\) and \(d\), optimizing over \(d\)-dimensional strategies is a
finite compact semialgebraic optimization problem.  The variables are the real
and imaginary parts of a unit state vector and of finitely many positive
semidefinite POVM matrices whose sums are the identity; the payoff is a
polynomial with rational coefficients when \(G\) has rational data.  By real
quantifier elimination, or equivalently by effective compact approximation
with certified semialgebraic feasibility, one can approximate this restricted
optimum sufficiently to distinguish the alternatives ``at least \(5/6\)''
and ``at most \(1/2\)''.

Thus \(D\) would yield an algorithm deciding the Halting Problem on the
instances encoded by the \(\mathrm{MIP}^{*}=\mathrm{RE}\) reduction, a
contradiction.  Hence no computable uniform dimension bound of the requested
kind exists.
\end{proof}

\section*{Scope and Interpretation}

This does not contradict the compactness assertion in the source paper.  For
each fixed pair \((\varepsilon,n)\), a finite bound exists non-effectively.
The theorem says that the requested explicit bound cannot be made computable.
It is therefore a negative answer to the algorithmic interpretation of
``provide an explicit bound'' in Open Problem 5.3.

This packet is also separate from Open Problem 5.2 of the same source paper,
which is already represented in the run by the packet recording Slofstra's and
Ji--Leung--Vidick's finite games requiring infinite-dimensional entanglement.

\section*{Verification Notes}

The local source paper is included as \texttt{source\_paper.pdf}.  The local
supporting paper is included as \texttt{supporting\_paper\_2001.04383.pdf}.
The relevant MIP* input is the abstract statement that the Halting Problem
reduces efficiently to deciding whether a two-player nonlocal game has
entangled value \(1\) or at most \(1/2\).

\begin{thebibliography}{9}

\bibitem{PVSurvey}
C. Palazuelos and T. Vidick,
\emph{Survey on Nonlocal Games and Operator Space Theory},
J. Math. Phys. 57 (2016), 015220;
arXiv:1512.00419.

\bibitem{MIPRE}
Z. Ji, A. Natarajan, T. Vidick, J. Wright, and H. Yuen,
\emph{\( \mathrm{MIP}^{*}=\mathrm{RE} \)},
Commun. ACM 64 (2021), no. 11, 131--138;
arXiv:2001.04383.

\end{thebibliography}

\end{document}
